\begin{center}
\textsc{Gauss} to \textsc{Olbers}.\quad Zeven, 8~July 1824.
\end{center}

\par\noindent\hspace*{\parindent}I duly received your last two kind letters .\ .\ .\ .\ .\ .

\par\noindent\hspace*{\parindent}The first has quite disheartened me about all my measurements.
By rejecting the Bremen-Brüttendorf-Wentel triangle \textls[120]{because the
angle at Bremen is too acute}, you thereby also condemn what seems to be
the only possible way of getting around Bremen on the other side, for in the
Bremen-Bottel-Steinberg triangle the angle at Bremen will be much more
acute still. You do add, with your usual kindness, that in the case of that
triangle the accuracy would not suffer appreciably
\par\clearpage
\noindent{\small\textbf{ON THE HANOVER TRIANGULATION.}}\hfill371\par\medskip
\noindent\textls[120]{on that account}, because I take such care to make my measurements precise.
But that reason, whose truth I shall leave open for now, can \textls[120]{not} reassure me
in the slightest.
My principle is that one should always observe as accurately as one can.
The degree of accuracy of the observations, no matter how high or low it
may be, always determines in turn the degree of accuracy one may require
of the results. And in my opinion the accuracy of the observations cannot
make an intrinsically bad triangle good. Otherwise it would at least have
been superfluous to measure the other good triangles with the same precision.
At best, such precision can then bring the whole system back on a par with
other intrinsically much poorer surveys.

\par\noindent\hspace*{\parindent}Up to now I have held it to be mere prejudice to consider triangles
with very small angles detrimental to \textls[120]{accuracy}, provided that the sides
opposite the acute angles do not serve as transition sides. I have considered
such triangles with small angles less satisfactory \textls[120]{only} because they
do not take one very far at a single step, and therefore require more time and
expense than if one can advance a long way in one step. Even this reason
falls away altogether if locating and preparing a large triangle takes perhaps
twice as long as measuring two triangles together that reach the same place,
one of which has a very acute angle. Even so, I have not acted entirely on
this principle. I have never adopted a triangle with a small angle until I had
exhausted \textls[120]{almost all} possibilities of avoiding it (except for the one that would
have led to poor measurements themselves, i.e.\ \textsc{Zach}'s high towers).
I had no expectation of thereby gaining anything in \textls[120]{accuracy}.
My motive was the perhaps pardonable wish to give the system as much beauty
and completeness of form as possible in addition to its intrinsic merit.

\par\noindent\hspace*{\parindent}But what I wrote about this in an earlier letter has failed to convince
\textls[120]{you}, and you require the precision of the measurements to make up for
the harm to accuracy that would otherwise arise from the acute angle, which
in my view is \textls[120]{impossible}. I am thus quite at a loss about my former view
\par\smallskip\noindent\hfill{\footnotesize 47\textsuperscript{*}}\par
\clearpage
\noindent372\hfill{\small\textbf{CORRESPONDENCE WITH OLBERS.}}\par\medskip
\noindent and doubtful whether it was not quite wrong, and at any rate I cannot
hope to convince others that it is right. As for the Bremen-Brüttendorf-Wentel
triangle in particular, I myself would have been very reluctant to choose it
because it is not \textls[120]{beautiful} and does not take one very far at one step.
As regards \textls[120]{accuracy}, however (quite apart from how accurate the angle
measurements themselves are), taking its angles as \(12^{0},46^{0},122^{0}\), I
would have put it \textls[120]{entirely} on a par with another whose angles were
\(76^{0},46^{0},58^{0}\).
